% ==================================================================
%  REMOVED DUPLICATES  (restructuring 2026-09-30)
%  This file is NOT \input in main.tex. It keeps every passage that
%  was dropped because the same content is kept elsewhere, together
%  with where it came from and where the kept version now lives.
%  Delete this file once you are happy with the restructuring.
% ==================================================================

% ------------------------------------------------------------------
% 1. From background.tex, "Wrenches, Contact Forces, and Internal
%    Forces" -> "Contact Forces and Wrenches".
%    Kept version: theory.tex, Sec. Grasp Matrix -> "Contact Forces
%    and Wrenches" (eq:lambda-stacked, eq:wrench-map), rewritten for
%    the planar case. Original (3D wrench, broken LaTeX):
% ------------------------------------------------------------------
% \begin{equation}
% \boldsymbol{\lambda}
% =
% \begin{bmatrix}
% \boldsymbol{\lambda}_1 \
% \boldsymbol{\lambda}*2 \
% \vdots \
% \boldsymbol{\lambda}*{n_c}
% \end{bmatrix}.
% \end{equation}
% The grasp matrix $G$ maps the contact forces to the resulting wrench
% acting on the object:  \mathbf{w} = G\boldsymbol{\lambda},
% where the wrench is given by  \mathbf{w} = [\mathbf{f}; \boldsymbol{\tau}] \in \mathbb{R}^{6}.

% ------------------------------------------------------------------
% 2. From background.tex, "Wrench Directions": the separate equation
%    rank(G) = 6 ("In three dimensions, this requires").
%    Kept version: theory.tex, Sec. Nullspace Cases, eq:rank-condition
%    (states n_nu = 6 in 3D and n_nu = 3 in the plane).
% ------------------------------------------------------------------

% ------------------------------------------------------------------
% 3. From background.tex, "Wrench Acting on the Object":
%    "The wrench acting on the object is the net effect of all contact
%     forces and moments. It is obtained through  w = G lambda."
%    and the boxed mapping  lambda --G--> w  with the sentence
%    "where lambda describes the contact forces and w describes their
%     combined effect on the object."
%    Kept version: theory.tex, eq:wrench-map. The two-forces example
%    paragraph from this subsection was kept.
% ------------------------------------------------------------------

% ------------------------------------------------------------------
% 4. From background.tex, "Relation to the Force-Closure Test" (end):
%    the second boxed  lambda --G--> w  mapping, the repeated
%    G lambda = 0 equation, the boxed case distinction
%       G lambda != 0  => contact forces produce a wrench on the object
%       G lambda  = 0  => contact forces are internal
%    and the closing sentence:
%    "The rank of G determines which wrench directions can be generated,
%     while the nullspace of G describes contact-force combinations that
%     do not generate a net wrench. Both properties are therefore
%     fundamental when analyzing force closure."
%    Kept version: theory.tex, Sec. Nullspace Cases (rank and N(G)
%    paragraphs say the same). Re-use the closing sentence as a summary
%    there if you like it.
% ------------------------------------------------------------------

% ------------------------------------------------------------------
% 5. From theory.tex, "Contact and Friction Model" (Handbook p. 678):
%    "Wrench intensities lambda in N(G) are referred to as internal
%     object forces. These wrenches are internal because they do not
%     contribute to the acceleration of the object, i.e. G lambda = 0.
%     Instead, these wrench intensities affect the tightness of the
%     grasp. Thus, internal wrench intensities play a fundamental role in
%     maintaining grasps that rely on friction."
%    Kept version: merged (paraphrased) into theory.tex, Sec. Nullspace
%    Cases -> "Internal forces" paragraph, with the p. 678 citation.
% ------------------------------------------------------------------

% ------------------------------------------------------------------
% 6. From margin_force_closure.tex, "Wrench Space and Force Closure":
%    "Every admissible contact force induces a wrench (f_x, f_y, tau) on
%     the robot body about a reference point P_0 (e.g. the body's center
%     of mass), with tau = (p_i - P_0) x f_i. Collecting the normal- and
%     tangential-force contributions of all contacts into the grasp
%     matrix G in R^{3 x 2n_c} (Sec. grasp-matrix) gives the linear map
%     from contact force coefficients c = [f_{n,1}, f_{t,1}, ...,
%     f_{n,n_c}, f_{t,n_c}]^T to the resulting body wrench:
%        w = G c,  w = (f_x, f_y, tau) in R^3."
%    Kept version: theory.tex, Sec. Grasp Matrix -> "Contact Forces and
%    Wrenches" (eq:wrench-map, same map with c renamed lambda).
% ------------------------------------------------------------------

% ------------------------------------------------------------------
% 7. From margin_force_closure.tex, "LP2" (first half):
%    "Following the frictional form-closure test of Bicchi et al.
%     (Springer Handbook of Robotics, Sec. 38.4.2), the friction cone at
%     each contact is linearized into a set of half-space constraints
%     F c <= 0, where F collects the facet normals of every contact's
%     linearized friction cone. ..."
%    \begin{equation}
%    \begin{aligned}
%    \underset{\mathbf{c},\, d}{\text{maximize}} \quad & d \\
%    \text{subject to} \quad & G\mathbf{c} = \mathbf{0}, \\
%    & F\mathbf{c} \le -\mathbf{1}\, d, \\
%    & d \ge 0, \\
%    & \mathbf{e}^\top \mathbf{c} \le n_c,
%    \end{aligned}
%    \end{equation}
%    "where 1 and e are vectors of ones of appropriate dimension, and
%     the normalization constraint e^T c <= n_c prevents the trivial
%     unbounded scaling of c."
%    Kept version: theory.tex, Sec. LP2 (eq:lp2), taken from
%    implementation.tex. NOTE the differences: sign convention of F
%    (F >= 0 matches eq:face-normal-full) and e (selects normal
%    components, not all ones). The intro sentence was kept.
% ------------------------------------------------------------------
